\begin{remark}[The data that complete the one-loop components]
\label{rem:c11-next}
The next computations are the following. On the action side: the arity-two equations along $(K_b,P_{b'})$ and $(\eta,\eta')$ at a finite endpoint, which contain the coefficients $\theta^K_d$ and $\kappa^K_j$ of Proposition \ref{prop:c11-oddsol} (v) and give the one-loop parts of $[K,P]=D+M$ and $\{S,S\}=K$, and along $(\xi_-,\eta_\pm)$; the three- and four-leg parts of the arity-one equations, which decide whether the sector parts of the generators extend to symmetries of the interacting theory, fix the differences of the sector coefficients (Proposition \ref{prop:c11-arity1} (iii), Proposition \ref{prop:c11-freexi} (ii)) and the couplings $\kappa_k$ of Theorem \ref{thm:c9-F3}; and the arity-three component of $s(\hbar F_1)=\hbar\,\mathrm{Ob}$. On the operator side: $Q_1$, $S_1$, $M_1$, $K_1$ and $P_1$ along generic directions on the families of Definition \ref{def:c11-families}; the arity-two operator kernels, with which the relations (R2)--(R6) of Proposition \ref{prop:c11-relations} become identities of local functionals before the projection to cohomology; and on the scalar bilinears the one-loop components on the level of Weyl weight $\tfrac52$, with which (R2) and (R3), the Konishi relation between $Q_1$, $S_1$ and $D_1=-48$ units, make the combinations of Corollary \ref{cor:c11-F0} (ii) scheme independent. The class-level comparison of Theorem \ref{thm:c11-twistcheck} beyond the chiral bilinears uses the transport map $\Phi$ built from [ESW, Thm 10.9] with the normalisation of the letters $b$, $\beta$, $\gamma$, $\partial_{\dot\alpha}c$ of [BK], the condition on $D$, the arity-two kernels along $(\xi,\xi)$, which give the entry $bc$ with $b$ represented by the self-dual field strength and the entry $bb$ on the pair $(A|_V,A|_V)$ with two derivatives, and the relations [BK, (3.59)--(3.63)] applied to the entries of Table \ref{tab:c11-twistentries}. The inner-product relations of Proposition \ref{prop:c11-S} (iv) and Corollary \ref{cor:c11-K} use $\mathrm{Gram}_{\rm 1}$ from the two-point graphs of Definition \ref{def:c8-2pt}.
\end{remark}
